\section{Hosts and the auxiliary characters}\label{sec:hosts}

\subsection{The host field}
Given a primitive character $\psi$, an integer $e_2$ and two primes $q$ and $d$, a \emph{host} is
a field $E$ of the following kind \lvS{} (\lean{IsChiField}): $E$ embeds in a cyclotomic field
$\Q(\zeta_N)$ with $2^{e_2+2}$, the conductor of $\psi$, $q$ and $d$ all dividing $N$, and the
automorphism $\zeta_N\mapsto\zeta_N^{r}$ fixes $E$ exactly when
\[
r\equiv\pm1\pmod{2^{e_2+2}},\qquad \psi(r)=1,\qquad \Bigl(\frac rq\Bigr)=1,\qquad
\Bigl(\frac rd\Bigr)=1 .
\]
So the Galois group of $E$ is the product of a group of order four, coming from $q$ and $d$,
and a group $H$ built from a layer of the cyclotomic $\Z_2$-tower and from $\psi$. It is the
field to which \cref{st:vwin} applies. The host contains $\sqrt{-d}$, and $\Phi$ is the subgroup
that fixes it (\lean{HostRaw}); so the CM field of \cref{sec:period,sec:divisibility} is
$\Q(\sqrt{-d})$, and $\chi_\Phi$ is its quadratic character.

The sizes are fixed by two parameters \lvS{} (\lean{WinPair}, \lean{PageN}): a parameter $Y$,
which is constant on a window of scales, and the scale $x$ itself. The prime
$q$ satisfies $q\equiv3\pmod8$ and $\tfrac12Y^{\eps/5}<\log q\le Y^{\eps/5}$; the prime $d$
satisfies $d\equiv7\pmod{16}$ and $\log d\approx Y^{\eps/5}(\log Y)^2$, and $d+1$ is divisible
by $q$ and by every odd prime up to $Y^{\eps/5}$; and $2^{e_2}$ lies between $x^{2+8\eps}$ and
$2x^{2+8\eps}$. The scales of a window lie between $Y$ and a fixed power of $Y$, so the
conductor of $E$ exceeds every power of $x$, while the degree of $E$ is a power of $x$.

\subsection{Supply}
\begin{statement}[\lean{supplyChi}]\label{st:supply}
For every primitive $\psi$, if the cohomological bound holds for $\psi$ with exponent $401$,
then hosts with the required properties exist at every scale in a set containing windows of
exponent $\tfrac32$.
\end{statement}
The only assumption about $\psi$ is its cohomological bound. The windows come from
\cref{st:vwin}, which holds on windows $[Y_1,Y_1^2]$; this is why \cref{st:landau} is needed in its
windowed form.

\subsection{The auxiliary characters and the absence of a circle}
The reduction of \cref{sec:reduction} assumed $\abs{G_\chi(x)}\le x^{1/2+\eps}$ at the point
$x$ for the characters of the host that are not kept. They are the twists
$\xi\cdot\bigl(\tfrac\cdot q\bigr)^i\bigl(\tfrac\cdot d\bigr)^j$ with $(i,j)\ne(0,0)$ of
characters $\xi$ whose modulus is that of $\psi$. That looks like the statement being proved.
It is supplied by the choice of the primes:
\begin{statement}[\lean{PairEx}]\label{st:pairs}
For every modulus $M_0$, every $0<\eps<\tfrac1{52}$ and every $A\ge1$, for all large $Y$ there
are primes $q$ and $d$ of the sizes above, prime to $M_0$, such that for every character $\xi$
modulo $M_0$ and every $(i,j)\neq(0,0)$ the twist $\chi$ of $\xi$ satisfies
\[
\abs{G_\chi(y)}\le y^{1/2+\eps/2}\quad\text{and}\quad
\Bigl|\frac{L'}{L}(1,\chi)\Bigr|\le4\log y\qquad\text{for all } Y\le y\le Y^{A},
\]
and such that the twists of the real characters of modulus dividing $8M_0$ have no zero with
$\re s>1-\tfrac1{400}$ up to a stated height.
\end{statement}
There is no circle, for a reason of size: the conductors of the twists are at least
$q>\exp(\tfrac12Y^{\eps/5})$, far larger than any power of $y$, so the bound concerns sums much
shorter than the conductor, and it is required only for one suitable pair $(q,d)$, not for a
given one. The archive proves the statement from a mean-square estimate over prime moduli
\lvS{} (\lean{meanSqK}) and a zero-density estimate \lvS{} (\lean{zero_density_box}) for zeros
with $\re s>1-\tfrac1{400}$, in the tradition of Bombieri \cite{bombieri1965}; compare
\cite{siegel1935,page1935} for the exceptional real zero that the second clause excludes.

\subsection{Two combinatorial devices}
The thick set of \cref{sec:reduction} is built from Behrend's construction
\cite{behrend1946}, which the archive takes from Mathlib \lvS{} (\lean{noRel_image},
\lean{hostThickExists}). The property used is stronger than the absence of three-term
progressions: for three distinct elements $s_1,s_2,s_3$ of the set and distinct non-zero
integers $i,j$ of absolute value at most a given bound, one has $j(s_2-s_1)\ne i(s_3-s_1)$.

The statement $\mathsf{Green}$ of \cref{sec:overview} is \lvS: for every prime $p$ and every
$\nu\ge1$, the cuspidal characters of $\GL_\nu(\F_p)$ are parametrised, without repetition, by the
monic irreducible polynomials of degree $\nu$ with non-zero constant term, in a fixed labelling,
and each has degree $\prod_{i=1}^{\nu-1}(p^i-1)$. This is part of Green's theorem
\cite{green1955}.
